\documentclass[a4paper,12pt]{article}
\usepackage[T1]{fontenc}
\usepackage[utf8]{inputenc}
\usepackage[italian]{babel}
\usepackage{amsmath,cases,eqparbox,amsthm}
\usepackage{geometry}
\geometry{a4paper,top=2cm,bottom=2cm,left=0.8cm,right=1.0cm,%
	heightrounded,bindingoffset=5mm}
\usepackage{graphicx}
\graphicspath{ {./images/} }

\usepackage{multicol}
\setlength{\columnsep}{1cm}

\usepackage{tasks}
\usepackage{exsheets}
\SetupExSheets[question]{type=exam}

\title{Verifica di {MAT} }
\author{crever}

\theoremstyle{plain}
\newtheorem{esercizio}{Esercizio}

\renewcommand{\thefootnote}{\fnsymbol{footnote}}
\pagenumbering{gobble}

\begin{document}
%\maketitle

\begin{center}
	{\textbf{Verifica di {MAT}} \verb!IDV! \textbf{- {IST} }}
\end{center}

\small{\textit{\textbf{ISTRUZIONI:} Gli esercizi a risposta multipla hanno una sola risposta esatta. Per ogni risposta corretta si ottengono punti 1 e per le risposte errate non si perdono punti. Gli esercizi contrassegnati con (*) rappresentano gli obiettivi minimi. La
valutazione di base è pari a 2 punti}}

\begin{center}
	\begin{tabular}{llll}
		\hline
		&  &  & \\
		Cognome .................................... & 
		Nome ....................................... &
		Classe ......... & 
		Data .................\\
		\hline
	\end{tabular}
\end{center}

% \bigskip

\begin{multicols}{2}
	
\settasks{
	counter-format=[tsk[A]],
	label-width=4ex
}

% begin template esercizio
%<<SECESR>>
\begin{esercizio}
	%<<SECTPL>>	
	\emph{\textbf{ASR} {CMD}} %{LVL}
	\medskip
	\newline
	\small{ESR}
	% \hspace*{\fill} {\verb!(punti PNT)!}
	\bigskip
	%<<SECTPL>>	
\end{esercizio}	
%<<SECESR>>
% end template esercizio

\end{multicols}

\end{document}
